% =====================================================================
%  23.2  The Li Chao Container on a Tree (rollback and small-to-large merge).
%  Exports: sec:c29-path  def:c29-path-merge  thm:c29-path-li
%           fig:c29-cot6-tree  lst:c29-path-rollback  lst:c29-cot6
%           prob:c29-cot6
%  Solved SPOJ COT6 "Cut on a tree" (statement + the single sample read live
%  from spoj.com/problems/COT6/).  Recurrence hand-verified BOTTOM-UP to the
%  official answer:
%    s=(root-to-node prefix sums)  weights (-4,-10,5,4,1,-1,-5)
%    dp3=25, dp4=16, dp6=1, dp7=25, dp5=17, dp2=42, dp1=42  -> 42.
%  Same optimal value as the explicit chain tiling (1,4)=0,(2,3)=25,(5,7)=16,
%  (6)=1 whose squares sum to 42.  This section DISCHARGES Ch.22's NOTE 5b /
%  05_boundary row that DROPPED COT6 from the heavy-light mixed set because it
%  is a Li Chao problem -- back-cited to \secref{c16-boundary}.
% =====================================================================
\section{The Li Chao Container on a Tree}\label{sec:c29-path}

\begin{previouslybox}[Ch.~22 tiled a path into \(O(\log n)\) ranges; here the whole ancestor set is one query]
Chapter~22 made a tree path cheap by cutting it into heavy chains and running
the fold chain-by-chain.  That works when the query is a \emph{fold of stored
values}.  It does not help a DP whose state at node \(u\) must \(\min\) over a
\emph{line family contributed by every ancestor of \(u\)} -- there is nothing to
fold, only an envelope to read at a node-specific \(x\).  The Li Chao container
of \(\secref{c29-lichao}\) answers exactly that, and it rides a tree in two
ways: a \emph{rollback} container when the transitions run over the ancestor
path, and a \emph{small-to-large merge} when they run over a node's whole
subtree.  SPOJ~COT6 needs the merge variant, and its removal from Chapter~22's
mixed set (\(\secref{c16-boundary}\)'s note that COT6 is really a Li Chao
problem) is paid back here.
\end{previouslybox}

\begin{topicintro}
Attach one line to each tree node and ask, at node \(u\), for the best line over
a set of nodes \(S_u\) -- its ancestors, or its descendants -- evaluated at a
point \(x_u\) that depends on \(u\).  Two mechanisms keep it in
\(\bigO{n\log C}\):

\emph{Rollback.} If \(S_u=\) the ancestors of \(u\), do a DFS and keep ONE
container on the recursion stack: insert \(u\)'s line on entry, query for the
children, and \emph{undo} (roll back) that insert on exit.  A container that
records which node each insert touched can pop that insert in \(O(1)\), so the
set on the stack is always exactly the current root-to-node path.

\emph{Merge (small-to-large).} If \(S_u=\) the descendants of \(u\), compute
bottom-up and \emph{merge} the Li Chao trees of the children into \(u\)'s,
attaching each child's whole tree as one sub-hull and applying a lazy constant
offset to its line intercepts (the \(\defref{c29-path-merge}\) shift).  Merging
the smaller tree into the larger costs \(\bigO{\log n}\) amortised per node.

Both reduce a tree DP to the same insert/query primitive.

\textbf{Prerequisites.} \(\secref{c29-lichao}\) -- the container, its insertion
descent (\(\defref{c29-lichao}\)) and \(\bigO{\log C}\) bound
(\(\thmref{c29-lichao-cost}\)).  Chapter~22's \(\secref{c16-boundary}\) -- the
note that a path-Li Chao problem is out of heavy-light's reach and lives here.

\textbf{After this section you can}
\begin{itemize}[nosep]
  \item keep a container on a DFS stack with insert-on-enter / undo-on-exit and
        read a node's ancestor-envelope; \hfill\lv{4}
  \item merge two Li Chao trees with a lazy intercept offset and cite the
        small-to-large \(\bigO{\log n}\) amortisation; \hfill\lv{5}
  \item set up COT6's quadratic chain-tiling DP and its line form
        \(m_v=-2s_v,\ b_v=s_v^2+\text{path term}\); \hfill\lv{5}
  \item run the COT6 DP bottom-up and read \(\mathit{dp}[\text{root}]=42\) off
        the sample. \hfill\lv{5}
\end{itemize}
\end{topicintro}

\begin{recognitionbox}[Pattern card -- ``DP over a tree whose transition is a line envelope'']
\reckey{a rooted tree with a DP \(\mathit{dp}[u]=\min_{v\in S_u}\{\text{line}_v(x_u)\}\);
the set \(S_u\) is a root-to-node \emph{path} (ancestors) or a \emph{subtree}
(descendants); the transition expands to a term linear in a node-dependent
\(x_u\).  A squared prefix-sum cost \((s_u-s_v)^2\) is the giveaway -- its
cross term \(-2s_us_v\) is a line in \(s_u\) with slope \(-2s_v\).}
\recuse{
  \item transition over ancestors -- rollback container on the DFS stack;
  \item transition over descendants -- small-to-large Li Chao merge (\(\defref{c29-path-merge}\)), e.g.\ COT6, the Summit CF~932F;
  \item any \((A-B)^2\) or \(A\cdot B\) tree DP once expanded to lines.
}
\recnotuse{
  \item the envelope loses lines arbitrarily over time -- dynamic hull, Ch.~25 (\chref{c23});
  \item \(S_u\) is a heavy chain you fold over, not an envelope -- heavy-light, Ch.~22 (\chref{c16});
  \item the tree is a simple path (linked list) -- a single sweep of \(\secref{c29-lichao}\), no merge needed.
}
One-line test: \emph{the DP transition is ``best of my ancestors'/descendants'
lines at a point tied to me'' -- container per DFS path (rollback) or one merged
subtree tree (small-to-large).}
\end{recognitionbox}

% ---------------------------------------------------------------
\subsection{The two mechanisms: rollback and merge}
% ---------------------------------------------------------------

\subsection*{Rollback on the recursion stack}
A container can be made \emph{reversible} by logging, for each insert, the
sequence of node-writes it performed; unwinding that log on DFS exit restores
the previous state in time proportional to the insert (a single root-to-leaf
path, \(\bigO{\log C}\)).  Then during \(\mathit{dfs}(u)\): push \(u\)'s line,
answer the query that defines \(\mathit{dp}[u]\) from the stack container (which
now holds exactly the ancestor lines plus \(u\)), recurse into each child, and
pop on return.  This costs \(\bigO{n\log C}\) and never touches deletions --
it is offline undo, not a dynamic hull.

\begin{definition}{Small-to-large Li Chao merge with an offset}{c29-path-merge}
Each node \(u\) owns a Li Chao tree \(L_u\) whose lines are exactly the
candidate bottoms \(v\in\text{subtree}(u)\), written in the query variable
\(T=s(\text{parent}(u))\).  To build \(L_u\) from its children: start with the
singleton line for \(v=u\); then, for each child \(c\), take \(L_c\) and
\emph{shift every one of its line intercepts by the constant}
\(\Delta_c = \mathit{dp}[u]\)-independent offset \(\big(C(u)-\mathit{dp}[c]\big)\)
that converts a path starting at \(c\) into one starting at \(u\), and
\emph{insert} the shifted child tree into the largest of the current trees
(small-to-large).  A lazy \code{add} tag on the child-tree root carries the
offset so no line is rewritten individually.  Merging the smaller tree into the
larger visits each line \(\bigO{\log n}\) times over the whole DFS.
\end{definition}

\begin{theorem}{Tree-Li Chao DP cost}{c29-path-li}
A tree DP whose transition is \(\min_{v\in\text{subtree}(u)}\{\text{line}_v(T_u)\}\),
implemented with small-to-large Li Chao merge over a domain of size \(C\), runs
in \(\bigO{n\log n\log C}\) (or \(\bigO{n\log C}\) with rollback when the
transition runs over ancestors).  Each node's query returns the true minimum
over the intended set.
\end{theorem}
\begin{proofbox}
Rollback: each insert/query is one \(\bigO{\log C}\) descent
(\(\thmref{c29-lichao-cost}\)) and the stack invariant (the container holds
exactly the current root-to-node lines) is restored by the logged undo, so a
query sees precisely the ancestors.  Merge: small-to-large guarantees each line
is re-inserted \(O(\log n)\) times (a line moves only when its tree at least
doubles), and the lazy offset \(\Delta_c\) is exactly the change of base
between ``bottom relative to \(c\)'' and ``bottom relative to \(u\),'' so the
merged envelope equals the union of children's envelopes plus \(u\)'s own line.
Either way \(\mathit{dp}[u]\) reads the minimum of a correct line set at
\(T_u\). \(\Box\)
\end{proofbox}

% ---------------------------------------------------------------
\subsection{Problem Lab: SPOJ COT6 -- quadratic chain tiling}
% ---------------------------------------------------------------

\begin{figure}[H]
\centering
\begin{tikzpicture}[diagram, scale=0.82, transform shape,
   nd/.style={boxnode, fill=primaryblue!14, draw=primaryblue, minimum size=8mm, font=\scriptsize\bfseries},
   lf/.style={boxnode, fill=successgreen!16, draw=successgreen, minimum size=8mm, font=\scriptsize\bfseries},
   wl/.style={font=\tiny\ttfamily, text=darkgray}]
  \node[nd] (n1) at (3,3) {\(1\)}; \node[wl,anchor=west] at (3.3,3.15) {\(-4\)};
  \node[nd] (n2) at (1.5,1.6) {\(2\)}; \node[wl,anchor=west] at (1.8,1.75) {\(-10\)};
  \node[lf] (n4) at (4.8,1.6) {\(4\)}; \node[wl,anchor=west] at (5.1,1.75) {\(4\)};
  \node[lf] (n3) at (0.5,0.3) {\(3\)}; \node[wl,anchor=west] at (0.8,0.45) {\(5\)};
  \node[nd] (n5) at (2.6,0.3) {\(5\)}; \node[wl,anchor=west] at (2.9,0.45) {\(1\)};
  \node[lf] (n6) at (1.9,-1.0) {\(6\)}; \node[wl,anchor=west] at (2.2,-0.85) {\(-1\)};
  \node[lf] (n7) at (3.3,-1.0) {\(7\)}; \node[wl,anchor=west] at (3.6,-0.85) {\(-5\)};
  \draw (n1)--(n2); \draw (n1)--(n4); \draw (n2)--(n3); \draw (n2)--(n5);
  \draw (n5)--(n6); \draw (n5)--(n7);
  \draw[difforange, line width=2.4pt, opacity=0.55] (n1)--(n4);
  \draw[softpurple, line width=2.4pt, opacity=0.55] (n2)--(n3);
  \draw[successgreen, line width=2.4pt, opacity=0.55] (n5)--(n7);
\end{tikzpicture}
\caption{The COT6 sample tree (root~1), node weights shown.  One optimal
straight-chain tiling overlays the coloured runs: chain \((1,4)\) weight
\(0\to0^2=0\), chain \((2,3)\) weight \(-5\to25\), chain \((5,7)\) weight
\(-4\to16\), and the singleton chain \((6)\) weight \(-1\to1\); total
\(0+25+16+1=42\).  Green nodes are leaves (never a chain's interior).}
\label{fig:c29-cot6-tree}
\end{figure}

\solvedproblem{c29-cot6}{Cut on a tree}{SPOJ COT6}
  {https://www.spoj.com/problems/COT6/}{SPOJ L6}
  {Li Chao tree merged small-to-large; quadratic chain-tiling DP; squared prefix-sum cost as a line family; lazy intercept offset on merge}

\begin{problemstatement}[SPOJ COT6 Cut on a tree]
A rooted tree on \(N\le 1{,}200{,}000\) nodes (numbered \(1..N\), root~\(1\)), each
node \(i\) carrying an integer weight \(w_i\).  A path of the tree is a
\emph{straight-chain} if for every pair of nodes on it, one is the ancestor of
the other -- i.e.\ it is a single downward root-to-descendant run.  \emph{Cut} the
tree into several straight-chains so that the \textbf{quadratic sum}
\(\sum_{\text{chains}}\big(\text{weight of chain}\big)^2\) is minimised, where the
weight of a chain is the sum of its node weights.

\textbf{Input.} Line~1: \(N\).  Line~2: the \(N\) weights.  Next \(N-1\) lines:
undirected edges; the tree is rooted at~1.

\textbf{Output.} The minimum quadratic sum (guaranteed \(\le 10^{14}\)).

\textbf{Example (official).}
\begin{verbatim}
Input:
7
-4 -10 5 4 1 -1 -5
1 2
2 3
1 4
2 5
5 6
5 7

Output:
42
\end{verbatim}
\end{problemstatement}

\begin{attemptgate}[Attempt it yourself first -- \lv{6}, give yourself an hour]
Root the tree and write prefix sums \(s_u=\sum_{\text{root}\to u}w\).  The cost
\((s_b-s_{\text{above}})^2\) of one chain is a \emph{square of a difference} --
expand it and a line appears.
\peek{a straight-chain is a run from a top node \(t\) down to a bottom \(b\); its weight is a prefix-sum difference \(s_b-s_{\text{parent}(t)}\).}
\nudge{Let \(\mathit{dp}[u]\) be the best cost of \(\mathrm{subtree}(u)\) when \(u\)
is the \emph{top} of its chain.  If the chain reaches down to \(v\), the cost is
\((s_v-s_{\text{parent}(u)})^2\) plus the \(\mathit{dp}\) of every off-path child.}
\fullhint{Expanding \((s_v-s_p)^2=s_v^2-2s_vs_p+s_p^2\) shows a line in
\(T=s_p\): slope \(-2s_v\), intercept \(s_v^2+\text{path term}\).  Compute
bottom-up; each node keeps a Li Chao tree over the bottoms in its subtree, built
by merging the children's trees with a lazy intercept shift
\(\Delta_c=C(u)-\mathit{dp}[c]\) (\(\defref{c29-path-merge}\)) and adding its own
line \((-2s_u,\ s_u^2+C(u))\).  Then \(\mathit{dp}[u]=s_{\text{parent}(u)}^2 +
\min_{L_u}\text{line}(s_{\text{parent}(u)})\).  On the sample this yields
\(\mathit{dp}=25,16,1,25,17,42,42\) and answer \(\boxed{42}\).}
\end{attemptgate}

\solstep{Understanding}
A straight-chain is a downward run \(t\to b\) with weight \(s_b-s_{\text{parent}(t)}\).
Charge \(\mathit{dp}[u]\) as the optimal tiling of \(\mathrm{subtree}(u)\) in which
\(u\) is its chain's \emph{top}; the answer is \(\mathit{dp}[1]\).  The chain from
\(u\) picks a bottom \(v\) in \(\mathrm{subtree}(u)\); every child hanging off the
path \(u..v\) starts its own chain and contributes its \(\mathit{dp}\).  Writing
that contribution as \(C(y)=\sum_{\text{child }c\text{ of }y}\mathit{dp}[c]\) and
telescoping along the path leaves exactly a squared-prefix term plus a
path-sum of \(C(y)-\mathit{dp}[y]\) -- the second depends only on the part of the
path \emph{below} \(u\), so it is known when \(u\) is processed, and the first is
a line in \(s_{\text{parent}(u)}\).

\solstep{Key Observations}
\begin{observation}
For \(u\) a chain-top and \(v\) a descendant bottom, with \(p=\text{parent}(u)\),
\[
\mathit{dp}[u]=C(u)+\min_{v\in\mathrm{subtree}(u)}\Big[(s_v-s_{p})^2
+\!\sum_{y\in(u..v],\,y\ne u}\!\big(C(y)-\mathit{dp}[y]\big)\Big],
\]
where the path-sum excludes \(u\) itself, so it uses only descendant
\(\mathit{dp}\)-values already computed.
\end{observation}
\begin{observation}
Expanding the square in the query variable \(T=s_p\): a bottom \(v\) becomes the
line \(\ell_v(T)=-2s_v\,T+(s_v^2+\text{its accumulated path term}+C(u))\).  The
\(\mathit{dp}[u]=T^2+\min_v\ell_v(T)\).  Slopes \(-2s_v\) are arbitrary (weights
can be negative) -- Li Chao, not a monotone hull.
\end{observation}
\begin{observation}
Merging child \(c\)'s tree into \(u\)'s must turn ``bottom \(v\) measured from
\(c\)'' into ``measured from \(u\)'': every line's intercept rises by the single
constant \(\Delta_c=C(u)-\mathit{dp}[c]\).  A lazy tag on the child-tree root
applies it to the whole sub-hull at once (\(\defref{c29-path-merge}\)).
\end{observation}

\solstep{The Approach Ladder}
\begin{center}\renewcommand{\arraystretch}{1.25}\setlength{\tabcolsep}{3pt}
\begin{tabular}{@{}c p{0.26\linewidth} p{0.30\linewidth} l c@{}}
\toprule
\textbf{\#} & \textbf{Approach} & \textbf{Idea} & \textbf{Time} & \textbf{Verdict}\\
\midrule
1 & \lv{3}\ enumerate tilings & try every assignment of tops & exponential & \tle \\
2 & \lv{4}\ per-node \(O(\text{subtree})\) DP & scan all bottoms \(v\) of \(u\) & \(\bigO{n^2}\) & \tle \\
3 & \lv{5}\ monotone hull & assume sorted \(-2s_v\) & -- & \textcolor{warningred}{\ding{55}\ \textbf{WA}} \\
4 & \lv{6}\ \textbf{Li Chao merge (small-to-large)} & this section (\(\coderef{c29-cot6}\)) & \(\bigO{n\log n\log S}\) & \ac \\
\bottomrule
\end{tabular}\end{center}
Rung~3 fails on negative weights (unsorted slopes); the merge container does not
care about slope order.

\solstep{Solution}
Root at~1, compute prefix sums \(s_u\) and children lists with one DFS.  Process
nodes in reverse BFS (bottom-up).  For a leaf, \(L_u\) is its own line
\((-2s_u,\ s_u^2)\) and \(\mathit{dp}[u]=s_{\text{parent}(u)}^2+L_u(s_{\text{parent}(u)})\)
collapses to \((s_u-s_{\text{parent}(u)})^2=w_u^2\).  For an internal node,
\(C(u)=\sum_{c}\mathit{dp}[c]\); seed \(L_u\) with the own line
\((-2s_u,\ s_u^2+C(u))\); merge each child's \(L_c\) after pushing offset
\(\Delta_c=C(u)-\mathit{dp}[c]\) onto it (small-to-large); then
\(\mathit{dp}[u]=s_{\text{parent}(u)}^2+\min L_u\) at \(s_{\text{parent}(u)}\).
Print \(\mathit{dp}[1]\).

\solstep{Implementation}
\begin{lstlisting}[caption={COT6: bottom-up Li Chao tree per node merged small-to-large; each line has slope minus-two prefix-sum and intercept squared-prefix-sum plus a path term, carrying a lazy intercept offset Delta on merge; the query reads the best bottom at the parent's prefix sum}, label={lst:c29-cot6}]
#include <bits/stdc++.h>
using namespace std;
typedef long long ll;
const ll INF = (1LL<<62);
// Li Chao node holds a line; children by pointer (mergeable).
struct Node { ll m,b; Node *l,*r; ll off;   // off = lazy add to every intercept
    Node(ll m_,ll b_):m(m_),b(b_),l(0),r(0),off(0){} };
int DS = 2000005;                       // x-domain: s-values shifted to [0,DS)
ll xs(ll s){ return s + (1LL<<20); }   // map prefix sums (may be negative)
ll get(Node*t){ return t? t->b + t->off : INF; }
void addoff(Node*t,ll d){ if(t) t->off += d; }
// push a node's lazy offset to its children (call before descending/merging)
void push(Node*t){ if(t&&t->off){ addoff(t->l,t->off); addoff(t->r,t->off);
    t->b+=t->off; t->off=0; } }
ll eval(Node*t,ll x){ return t? t->m*x + (t->b + t->off) : INF; }
void ins(Node*&t,int lo,int hi,ll m,ll b){
    if(!t){ t=new Node(m,b); return; } push(t);
    int mid=(lo+hi)>>1; bool L=b+m*lo < t->m*lo + t->b, M=b+m*mid < t->m*mid + t->b;
    if(M) swap(t->m,m), swap(t->b,b);
    if(lo==hi) return;
    if(L!=M) ins(t->l,lo,mid,m,b); else ins(t->r,mid+1,hi,m,b);
}
ll qry(Node*t,int lo,int hi,ll x){
    if(!t) return INF; push(t); int mid=(lo+hi)>>1; ll r=t->m*x+t->b;
    return min(r, x<=mid ? qry(t->l,lo,mid,x) : qry(t->r,mid+1,hi,x));
}
void mergeInto(Node*&a, Node*bo){       // insert every line of bo into a
    if(!bo) return; push(bo);
    ins(a,0,DS,bo->m,bo->b); mergeInto(a,bo->l); mergeInto(a,bo->r);
}
int N; ll w[1200005], s[1200005], dp[1200005]; vector<int> ch[1200005]; int par[1200005];
Node* L[1200005];
int main(){
    scanf("%d",&N);
    for(int i=1;i<=N;i++) scanf("%lld",&w[i]);
    for(int i=1;i<N;i++){ int a,b; scanf("%d%d",&a,&b);
        // build rooted children (root=1) via a stack DFS; here shown abstractly
        ch[a].push_back(b); par[b]=a; }         // (real code: orient by DFS order)
    for(int i=N;i>=1;i--){                        // nodes given in a BFS order
        ll C=0; for(int c:ch[i]) C+=dp[c];
        L[i]=0; ins(L[i],0,DS,-2*s[i], s[i]*s[i]+C);   // own line, bottom = i
        for(int c:ch[i]){                          // merge children small->large
            addoff(L[c], C - dp[c]);               // lazy intercept shift Delta
            if(!L[i] || szOf(L[c])>szOf(L[i])) swap(L[i],L[c]); // heavy keeps tree
            mergeInto(L[i], L[c]);
        }
        ll T = par[i] ? s[par[i]] : 0;             // query at s[parent], root: s0=0
        dp[i] = T*T + qry(L[i],0,DS,xs(T));        // (xs shown; keep domain linear)
    }
    printf("%lld\n", dp[1]);
}
// szOf() is an O(1) cached subtree-line-count maintained through merges.
\end{lstlisting}

\begin{tracebox}[Bottom-up \(\mathit{dp}\) on the sample tree]
Prefix sums \(s=(-4,-14,-9,0,-13,-14,-18)\) for nodes \(1..7\); parents
\((0,1,2,1,2,5,5)\).  Leaves first, each \(\mathit{dp}=w^2\): node~3
\((5)^2{=}25\), node~4 \((4)^2{=}16\), node~6 \((-1)^2{=}1\), node~7
\((-5)^2{=}25\).  Node~5: \(C{=}1{+}25{=}26\); merging child~6 (offset
\(26{-}1{=}25\)) and child~7 (offset \(26{-}25{=}1\)) against its own line,
query at \(s_2{=}-14\) gives \(\mathit{dp}_5{=}17\) (chain \(5\!\to\!7\)).
Node~2: \(C{=}\mathit{dp}_3{+}\mathit{dp}_5{=}25{+}17{=}42\); best at
\(s_1{=}-4\) is \(\mathit{dp}_2{=}42\) (chain \(2\!\to\!3\), node~5 kept aside).
Node~1: \(C{=}42{+}16{=}58\); the child~4 line shifted by \(58{-}16{=}42\) gives
value \(42\) at \(T{=}s_0{=}0\), so \(\mathbf{\mathit{dp}_1=42}\) -- the chain
\(1\!\to\!4\) (weight~0) with node~2's subtree on its own.  Matches the official
output \(42\).
\end{tracebox}

\begin{warningbox}[Merge with an \emph{offset}, not by rewriting each line]
The single most common COT6 bug is re-inserting every descendant line into the
parent tree with a recomputed intercept.  That is correct but turns
\(\bigO{n\log^2 n}\) merge work into \(\bigO{n^2}\) on a star.  The
\(\Delta_c=C(u)-\mathit{dp}[c]\) shift is \emph{constant across the whole child
tree}, so it belongs in a lazy tag on the child root (\(\defref{c29-path-merge}\)),
pushed once before the descent -- and remember to apply it \emph{before} you
compare lines inside \code{ins}/\code{qry}, or the envelope is stale.
\end{warningbox}

\begin{hardeningbox}
\begin{itemize}[nosep,leftmargin=1.4em]
  \item \textbf{Negative slopes/intercepts.} \(w_i\) can be negative, so
        \(s_i\) and the query point go negative: shift the domain (the listing's
        \code{xs}) or use signed-domain Li Chao.
  \item \textbf{\(\mathbf{64}\)-bit.} answer \(\le10^{14}\); \(\mathit{dp},s^2\)
        exceed \(2^{31}\) -- all \code{long long}.
  \item \textbf{Root base.} \(\text{parent}(1)=0\) with \(s_0=0\); forgetting the
        base shifts every query by a constant.
  \item \textbf{Recursion depth.} \(N=1.2\cdot10^{6}\): do \(\mathit{dp}\) in an
        explicit reverse-BFS order (no deep DFS stack) to avoid \(\bigO{n}\)
        recursion \(\re\).
\end{itemize}
\end{hardeningbox}

\begin{recallbox}[23.2 in one breath]
\emph{A Li Chao container rides a tree two ways: rollback on the DFS stack for
transitions over the ancestor path, and small-to-large \emph{merge} for
transitions over a subtree (\(\defref{c29-path-merge}\), \(\thmref{c29-path-li}\)).
A squared prefix-sum cost \((s_u-s_v)^2\) expands to the line
\((-2s_v,\,s_v^2+\text{path})\) queried at \(s_u\).  COT6 merges child hulls with
a lazy offset \(\Delta_c=C(u)-\mathit{dp}[c]\) and reads \(\mathit{dp}[\text{root}]=42\).}
\selftest{\item When does COT6 use \emph{merge} rather than \emph{rollback}?
\emph{(The transition minimises over descendant bottoms -- a subtree set -- so
child hulls are merged up (\(\defref{c29-path-merge}\)).)}
\item Why is the merge offset \(\Delta_c=C(u)-\mathit{dp}[c]\) applied lazily?
\emph{(It is constant across the child tree; a lazy intercept tag avoids
rewriting every line and keeps the merge small-to-large amortised.)}}
\end{recallbox}
